\chapter{Обзор методов дифференцируемого рендеринга для реконструкции и оптимизации конфигурации освещения}

Задача реконструкции конфигурации освещения по фотографиям, находящаяся на стыке компьютернрой графики и зрения, имеет бесконечное множество вариантов решения, поскольку изображение получается путём комбинации параметров источника освещения и материалов объекта, и при попытке их разделения возникают трудности. В связи с этим, большинство методов совместно реконструирует материалы и освещение сцены. В 2019 году был предложен обратный трассировщик путей \cite{Azinovic_2019_CVPR}, метод совместной оценки материалов объектов и источников освещения в сценах внутри помещений с использованием дифференцируемой трассировки путей методом Монте-Карло. Входными данными служат грубая геометрия сцены, позиции камер и соответствующие изображения; оптимизация проводится методом стохастического градиентного спуска по параметрам рассеивающей BRDF и положениям/интенсивностям источников освещения одновременно. Работа продемонстрировала, что корректное моделирование многократных отражений принципиально необходимо для разрешения неоднозначности между диффузным альбедо и прямым освещением. Однако трассировка путей для оптимизации является вычислительно затратной, из-за чего реконструкция сцены занимала несколько часов. Эта проблема решалась в методе 2023 года, факторизованной обратной трассировке путей \cite{fipt2023}. Ключевая идея состоит в разбиении интеграла рендеринга на произведение компонент BRDF и затенения: последние предварительно вычисляются и кешируются, а параметры BRDF и излучение источников оптимизируются независимо. Алгоритм обнаруживает источники и сходится существенно быстрее, лучше разрешая неоднозначность между параметрами материалов и источников. Метод допускает оптимизацию параметров источников менее чем за час на сценах в помещениях. Метод PhySG \cite{physg2021} одновременно восстанавливает геометрию, материалы и освещение сцены в помещении. Геометрия представлена функцией расстояния со знаком, хранящейся в многослойных перцептронах, а освещение и BRDF -- сферическими гауссианами. В 2020 году \cite{li2020inverse} был предложен глубокий фреймворк обратного рендеринга для сцен внутри помещений, восстанавливающий форму, пространственно-переменное освещение и SVBRDF из одиночного изображения. Освещение представлено как набор сферических гауссиан на воксельной сетке; дифференцируемый слой рендеринга совместно оптимизирует все три компоненты. MAIR \cite{Choi_2023_CVPR} в 2023 году обобщил предыдущий подход на случай нескольких входных изображений, предложив механизм внимания для агрегации информации с нескольких ракурсов и более точную пространственно-переменную трёхмерную карту освещения. Расширенный датасет OpenRooms позволяет обучать сеть на синтетических HDR-данных с надёжным обобщением на реальные сцены.

Отдельная группа методов использует возможности полей освещённости для реконструкции освещения. TensoIR \cite{Jin2023TensoIR} расширил тензорное разложение полей освещённости (TensoRF \cite{Chen2022ECCV}) до задачи обратного рендеринга, обеспечив совместную реконструкцию геометрии, отражательных свойств поверхности и внешнего освещения из нескольких изображений при неизвестных условиях освещения. Тензорное представление низкого ранга позволяет эффективно моделировать вторичное затенение и обрабатывать изображения, захваченные при нескольких различных условиях освещения одновременно. Метод NeILF \cite{yao2022neilf} использует нейронное поле падающего света -- 5D-функцию, параметризующую освещение в каждой точке пространства. В отличие от методов, предполагающих единственную карту окружения, NeILF моделирует пространственно-переменное освещение, включая межотражения между поверхностями. Ключевое ограничение, связывающее падающий и исходящий свет через уравнение рендеринга, делает возможным физически корректное разделение геометрии, материалов и освещения. Он получил развитие в виде метода NeILF++ \cite{zhang2023neilfpp}, использущего комбинацию нейронного поля падающего света и нейронного поля освещённости, которые связаны уравнением рендеринга. NVDiffRec \cite{Munkberg_2022_CVPR} продемонстрировал эффективный метод совместной оптимизации топологии, материалов и освещения из нескольких изображений. Система выводит треугольные меши с пространственно-переменными материалами и HDR-картой окружения, непосредственно совместимые с традиционными системами рендеринга. Для эффективного восстановления освещения введена дифференцируемая формулировка split-sum аппроксимации для интегрирования по окружению; дальнейшее развитие этой работы -- NVDiffRecMC \cite{nvdiffrecmc2022} -- заменяет приближённое интегрирование точным методом Монте-Карло с шумоподавлением.

Точное моделирование многократных отражений остаётся открытой проблемой. I$^{2}$-SDF \cite{zhu2023i2} и родственные методы решают её путём раздельного моделирования нейронной функции расстояния со знаком, поля яркости и поля излучения с последующей трассировкой путей для нейронного представления. TensoFlow \cite{gu2025tensoflow} вводит обучаемый пространственно- и направленно-осведомлённый сэмплер на основе тензорных нормализующих потоков для интеграла рендеринга: параметризованный тензорным представлением сэмплер точнее аппроксимирует подынтегральное выражение в сравнении с предопределёнными схемами выборки.

Для реконструкции освещения издалека, при отсутствии источников света внутри сцены, используются HDR-карты окружения. В 2021 году была предложена сеть EnvMapNet \cite{HDR_EM_estim} для оценки HDR-карт окружения из LDR-изображения с небольшим углом обзора в режиме реального времени (менее 9 мс на iPhone XS). Две новых функции потерь -- ProjectionLoss для генерируемого изображения и ClusterLoss для состязательного обучения -- позволяют сократить ошибку направления оцениваемых источников более чем на 50\% по сравнению с предшествующими методами.

Отдельным классом задач идёт оптимизация расположения светильников для достижения заданного распределения света. Это задача прямого проектирования, которая также решаема инструментами дифференцируемого рендеринга. В 2024 году \cite{ALT2024} был представлен интерактивный независимый от ракурсов фреймворк обратного рендеринга для непрерывной оптимизации расположения источников света в трёхмерной сцене. Ключевым отличием от большинства методов является использование трассировки от источников света, которая более естественна для задач, параметры оптимизации которых связаны с источниками. Метод вводит аналитическую сопряжённую формулировку для вычисления градиента целевой функции относительно параметров расположения светильников.

Быстрое развитие 3D Gaussian Splatting (3DGS) как представления сцены привело к появлению группы методов обратного рендеринга для реконструкции освещения с высокой скоростью вычислений. Метод relightable 3d gaussians \cite{R3DG2023} добавил к гауссианам новые свойства: нормали, параметры материалов и BRDF, а также падающий на них со всех направлений свет; реконструкция геометрии проводится с помощью сплаттинга гауссиан, а освещение -- через физически-обоснованный рендеринг. GS-IR \cite{GSIR2024} расширил 3DGS до задачи обратного рендеринга, оценивая геометрию сцены, свойства поверхности и освещение окружения на основе нескольких изображений, при этом условия освещения неизвестны. Для нормальной оценки вводится регуляризация на основе производных глубины, а непрямое освещение моделируется через запечённую карту окклюзии. Время обучения GS-IR примерно в 5 раз меньше по сравнению с аналогичными методами, основанными на NeRF, при сопоставимом качестве реконструкции. SVG-IR \cite{SVGIR2025} расширил постоянные параметры материала гауссиан до пространственно-переменных, что позволяет разным частям одной гауссианы иметь различные нормали и свойства материала. Интеграция физической модели непрямого освещения обеспечивает прирост PSNR на 3,5 дБ в сравнении с существующими методами, основанными на гауссианах.




\endinput